% ભાગ: first-order-logic
% પ્રકરણ: beyond
% વિભાગ: other-logics

\documentclass[../../../include/open-logic-section]{subfiles}

\begin{document}

\olfileid{fol}{byd}{oth}

\olsection{બીજાં તર્કશાસ્ત્રો}

અત્યાર સુધીમાં તમે સમજ્યા હશો કે નવું તાર્કિક તંત્ર રચવું મુશ્કેલ નથી.
તમે પણ તમારી પોતાની વાક્યરચના ઘડી, નિષ્પત્તિ-તંત્ર બનાવી અને તેને
અનુરૂપ અર્થવિજ્ઞાન તૈયાર કરી શકો. !!{derivation} તંત્રને અર્થવિજ્ઞાન
માટે પૂર્ણ બનાવવું હોય, તો થોડી ચતુરાઈ કરવી પડી શકે; અને તમારું
તર્કશાસ્ત્ર ખરેખર રસપ્રદ છે એવું વિશાળ જગતને મનાવવામાં થોડો પ્રયત્ન
લાગી શકે. પરંતુ બદલામાં, તમારા તર્કશાસ્ત્રના ગાણિતિક અને ગણનાત્મક
ગુણધર્મો શોધવામાં નિર્દોષ આનંદના અનેક કલાકો માણી શકો.

તાજેતરના દાયકાઓમાં વિધિવત્ તાર્કિક તંત્રોનો જાણે વિસ્ફોટ થયો છે.
અસ્પષ્ટ તર્કશાસ્ત્ર ધૂંધળા ગુણધર્મો વિશેના તર્કવિચારનું નિદર્શન કરવા
રચાયું છે. સંભાવ્યતા તર્કશાસ્ત્ર અનિશ્ચિતતા વિશેના તર્કવિચારનું
નિદર્શન કરવા રચાયું છે. ડિફૉલ્ટ અને અ-એકદિશવર્ધી તર્કશાસ્ત્રો પછીથી
નવી માહિતી સામે પાછાં ખેંચી શકાય એવા પરાજેય તર્કવિચારનાં સ્વરૂપો,
એટલે ``વાજબી'' અનુમાનોનું નિદર્શન કરવા રચાયાં છે. જ્ઞાન વિશેના
તર્કવિચારનું નિદર્શન કરતાં જ્ઞાનમીમાંસક તર્કશાસ્ત્રો; કારણાત્મક
સંબંધો વિશેના તર્કવિચારનું નિદર્શન કરતાં કારણાત્મક તર્કશાસ્ત્રો; અને
નૈતિક તથા આચારલક્ષી કર્તવ્યો વિશેના તર્કવિચારનું નિદર્શન કરતાં
``કર્તવ્યલક્ષી'' તર્કશાસ્ત્રો પણ છે. આ તંત્રો રજૂ કરવાની મુખ્ય પ્રેરણા
તત્ત્વચિંતનાત્મક, ગાણિતિક કે ગણનાત્મક છે તેના આધારે, આવા જીવોનો અભ્યાસ
ગાણિતિક તર્કશાસ્ત્ર, તત્ત્વચિંતનાત્મક તર્કશાસ્ત્ર, કૃત્રિમ બુદ્ધિ,
સંજ્ઞાનવિજ્ઞાન અથવા બીજા કોઈ શીર્ષક હેઠળ થતો જોવા મળશે.

યાદી લંબાતી જ જાય છે અને શક્યતાઓ અનંત લાગે છે. સમગ્ર માનવીય
તર્કબુદ્ધિને ગણતરીમાં ઘટાડવાનું લાઇબ્નિત્સનું સ્વપ્ન કદાચ આપણે ક્યારેય
સિદ્ધ ન કરીએ---પણ તેથી પ્રયત્ન કરવાનું અટકવાનું નથી.

\end{document}
